\documentclass[11pt]{article}

\usepackage[margin=1in]{geometry}
\usepackage{amsmath,amssymb,amsthm}
\usepackage{graphicx}
\usepackage[hidelinks]{hyperref}

\newtheorem{theorem}{Theorem}
\newtheorem{lemma}{Lemma}
\newtheorem{remark}{Remark}

\newcommand{\R}{\mathbb{R}}
\newcommand{\SX}{S_X}
\newcommand{\Iso}{\mathrm{Iso}(X)}
\newcommand{\dist}{\operatorname{dist}}
\newcommand{\Lin}{\mathcal{L}(X)}
\newcommand{\norm}[1]{\left\lVert #1 \right\rVert}

\title{Smooth finite-dimensional Banach spaces satisfy Property (P)}
\author{fa\_banach\_001, lane 19}
\date{2026-06-19}

\begin{document}
\maketitle

\section*{Source question}

This packet concerns arXiv:2407.07490, Debmalya Sain, Arpita Mal, Kalidas
Mandal and Kallol Paul, \emph{On uniform Bishop-Phelps-Bollobas type
approximations of linear operators and preservation of geometric properties}.
In Section 3 the authors define Property (P) and ask:

\begin{quote}
Is it true that each finite-dimensional Banach space \(X\) satisfies Property
\((P)\)? What about the infinite-dimensional case?
\end{quote}

\begin{center}
\includegraphics[width=0.92\linewidth]{figures/open_problem_crop.png}
\end{center}

The source paper proves Property (P) for finite-dimensional polyhedral spaces,
for Hilbert spaces, and later for \(\ell_p^2\) when \(p>2\) is an integer.  The
result below gives a new partial answer for all finite-dimensional smooth real
Banach spaces.

\section*{Definitions}

Let \(X\) be a real Banach space.  For \(A\in \Lin\), set
\[
  M_A=\{x\in \SX:\norm{Ax}=\norm{A}\}.
\]
The space \(X\) has \emph{Property (P)} if there is an \(r_0>0\) such that for
every norm-one non-isometry \(A\in \Lin\setminus \Iso\), there is
\(x_A\in \SX\) with
\[
  B(x_A,r_0)\cap M_A=\varnothing .
\]
Equivalently, the norm attainment set of a norm-one non-isometry can never be
uniformly dense in the unit sphere.

Recall also that a finite-dimensional real Banach space is smooth when every
\(x\in \SX\) has a unique norming functional \(j(x)\in S_{X^*}\).  In finite
dimension this implies that the norm is Frechet differentiable on
\(X\setminus\{0\}\), with derivative at \(x\in \SX\) equal to \(j(x)\), and the
map \(x\mapsto j(x)\) is continuous on \(\SX\).

\section*{Proof intuition}

Suppose Property (P) fails.  Then there are norm-one non-isometries
\(A_n\) whose norm-attainment sets \(M_{A_n}\) are \(1/n\)-dense in \(\SX\).
Such an operator almost preserves the norm of every unit vector, so compactness
forces \(A_n\) close to the compact isometry group.

Conjugate \(A_n\) by a nearest isometry and write the result as
\[
  B_n=I+t_n C_n,\qquad t_n=\norm{B_n-I}>0 .
\]
After passing to a subsequence, \(C_n\to C\) with \(\norm{C}=1\).  Contractivity
of \(B_n\) gives \(j(x)(Cx)\le 0\).  The density of \(M_{B_n}\) gives the
opposite inequality, hence
\[
  j(x)(Cx)=0\qquad(x\in \SX).
\]
That identity says \(C\) is an infinitesimal isometry: \(\exp(sC)\) is an
isometry for every \(s\in\R\).  But then the isometry \(\exp(t_n C)\) is
closer to \(B_n=I+t_n C_n\) than \(I\) is, contradicting the choice of the
nearest isometry.

\section*{Main theorem}

\begin{theorem}
Every finite-dimensional smooth real Banach space satisfies Property \((P)\).
\end{theorem}

\begin{proof}
Let \(X\) be finite-dimensional and smooth.  Denote by \(G=\Iso\) the group of
surjective linear isometries of \(X\).  Since \(X\) is finite-dimensional,
\(G\) is a compact subset of \(\Lin\).

Assume, toward a contradiction, that \(X\) fails Property \((P)\).  Then for
each \(n\) there exists a norm-one non-isometry \(A_n\in \Lin\) such that
\(M_{A_n}\) is \(1/n\)-dense in \(\SX\).

We first note that \(\dist(A_n,G)\to 0\).  Indeed, for every fixed
\(x\in \SX\) choose \(y_n\in M_{A_n}\) with \(\norm{x-y_n}\le 1/n\).  Since
\(\norm{A_n}=1\),
\[
  1\ge \norm{A_nx}\ge \norm{A_ny_n}-\norm{A_n}\norm{x-y_n}
    \ge 1-\frac1n .
\]
Thus any operator-norm cluster point \(A\) of the sequence \(A_n\) satisfies
\(\norm{Ax}=1\) for all \(x\in \SX\), hence belongs to \(G\).  Compactness of
the unit ball of \(\Lin\) proves \(\dist(A_n,G)\to 0\).

For each \(n\), choose \(U_n\in G\) with
\(\norm{A_n-U_n}=\dist(A_n,G)\), and put
\[
  B_n=U_n^{-1}A_n .
\]
Then \(B_n\) is a norm-one non-isometry, \(M_{B_n}=M_{A_n}\), and
\[
  \dist(B_n,G)=\norm{B_n-I}.
\]
The last equality follows because left multiplication by \(U_n^{-1}\) is an
isometry of \(\Lin\) and maps \(G\) onto \(G\).

Set
\[
  t_n=\norm{B_n-I}>0,\qquad C_n=\frac{B_n-I}{t_n}.
\]
After passing to a subsequence, \(C_n\to C\) in \(\Lin\), with
\(\norm{C}=1\).  We now show that
\[
  j(x)(Cx)=0\qquad (x\in \SX).
\]

First fix \(x\in \SX\).  Since \(\norm{B_n}\le 1\),
\[
  \norm{x+t_n C_nx}=\norm{B_nx}\le 1=\norm{x}.
\]
Dividing by \(t_n\) and using Frechet differentiability of the norm at \(x\),
with \(C_nx\to Cx\), gives
\[
  j(x)(Cx)\le 0. \tag{1}
\]

For the reverse inequality, choose \(y_n\in M_{B_n}\) with
\(\norm{y_n-x}\le 1/n\).  Then \(y_n\to x\), and
\[
  \norm{y_n+t_n C_ny_n}=\norm{B_ny_n}=1=\norm{y_n}.
\]
The norm is differentiable uniformly on compact subsets of \(X\setminus\{0\}\)
in finite dimension, and \(C_ny_n\to Cx\).  Therefore
\[
  0
  =\lim_n
    \frac{\norm{y_n+t_n C_ny_n}-\norm{y_n}}{t_n}
  =j(x)(Cx). \tag{2}
\]
Equations (1) and (2) give the desired identity.

We claim that \(\exp(sC)\in G\) for every \(s\in\R\).  Fix \(x\ne 0\) and set
\(g(s)=\norm{\exp(sC)x}\).  The vector \(\exp(sC)x\) is never zero.  If
\(v_s=\exp(sC)x/\norm{\exp(sC)x}\), then differentiability of the norm and the
identity \(j(v)(Cv)=0\) on \(\SX\) give
\[
  g'(s)
   = j(v_s)\bigl(C\exp(sC)x\bigr)
   = \norm{\exp(sC)x}\,j(v_s)(Cv_s)
   =0 .
\]
Hence \(g(s)=g(0)=\norm{x}\) for all \(s\).  Thus \(\exp(sC)\) is a linear
isometry of \(X\), and in finite dimension it is surjective, so \(\exp(sC)\in
G\).

Finally, since \(\exp(t_n C)\in G\),
\[
  t_n=\dist(B_n,G)\le \norm{B_n-\exp(t_n C)}.
\]
But \(B_n=I+t_n C_n\), \(C_n\to C\), and
\[
  \exp(t_n C)=I+t_n C+O(t_n^2),
\]
so
\[
  \norm{B_n-\exp(t_n C)}
  \le t_n\norm{C_n-C}+O(t_n^2)=o(t_n).
\]
This contradicts \(t_n>0\).  Therefore \(X\) must satisfy Property \((P)\).
\end{proof}

\section*{Scope and limitations}

This is a partial answer to the source question.  It proves the
finite-dimensional smooth real case.  It does not settle arbitrary
finite-dimensional norms, because nonsmooth nonpolyhedral spaces do not have a
global single-valued differentiable duality map.  It also does not settle the
infinite-dimensional part of the source question; the proof uses compactness
of both the operator unit ball and the isometry group.

The argument should be adaptable to complex spaces if smoothness is formulated
as real Frechet differentiability of the norm and \(j(x)\) is interpreted as
the real derivative.  This packet records only the real version to avoid hiding
that convention.

\section*{Verification notes}

\begin{itemize}
\item The failure of Property (P) is used exactly as the existence of
  norm-one non-isometries whose \(M_A\)'s are arbitrarily dense in \(\SX\).
\item The passage from dense \(M_{A_n}\) to closeness to the isometry group
  uses only finite-dimensional compactness.
\item The nearest-isometry conjugation is essential: it gives
  \(\dist(B_n,G)=\norm{B_n-I}\), so the later exponential isometry contradicts
  nearest-point minimality.
\item The only differentiability input is standard for finite-dimensional
  smooth real Banach spaces: the norm is Frechet differentiable off the origin,
  with continuous derivative on the sphere.
\item A bounded novelty check was performed against the run indexes for
  arXiv:2407.07490, Property (P), uniform BPB/Bishop-Phelps-Bollobas keywords,
  and nearby BPB phrases.  No duplicate packet or exact later answer was found.
\end{itemize}

\section*{Human review recommendation}

Review as a likely valid partial result.  The main line to audit is the
standard differentiability fact for finite-dimensional smooth real Banach
spaces and its use with \(y_n\to x\) and \(C_ny_n\to Cx\).  Once that lemma is
accepted, the compact Lie-tangent contradiction is quite rigid.

\section*{Reference}

D. Sain, A. Mal, K. Mandal and K. Paul, \emph{On uniform
Bishop-Phelps-Bollobas type approximations of linear operators and preservation
of geometric properties}, arXiv:2407.07490.

\end{document}

